% =====================================================================
%  1bat-ud3-4.tex  —  HORA 4: Funcions lineals: l'equació de la recta
%  (sense preàmbul: s'inclou des de main.tex)
% =====================================================================
\horatitol{Sessió 4 — Funcions lineals: l'equació de la recta}%
{Llegir les expressions de primer grau com a funcions: el pendent $m$ i el terme independent $n$, i el seu sentit social.}

\subsection{Idea clau}

A la sessió anterior trobàvem talls amb els eixos d'expressions com $-20x+600$.
Aquelles expressions de primer grau, que ja sabíem construir a la UD2, són
\textbf{funcions lineals}: la seva gràfica és una \textbf{recta}. Tota recta es pot
escriure com $y=mx+n$, i els dos números $m$ i $n$ ens expliquen tota la història.

\begin{idea}
\textbf{Equació de la recta:} $\quad y=mx+n$
\begin{itemize}[leftmargin=1.4em, itemsep=2pt]
  \item \textbf{$m$ és el pendent}: el \emph{ritme de creixement}. Indica quant puja
        (o baixa) $y$ cada cop que $x$ augmenta en $1$.
        Si $m>0$ la recta \textbf{creix}; si $m<0$ \textbf{decreix}; si $m=0$ és
        \textbf{horitzontal}.
  \item \textbf{$n$ és el terme independent} (l'\emph{ordenada a l'origen}): el valor
        de $y$ quan $x=0$. És el \textbf{punt de partida} i coincideix amb el tall amb
        l'eix vertical, $(0,n)$.
\end{itemize}
\textbf{Lectura social:} en una tarifa $C=mx+n$, el terme $n$ és la
\emph{quota fixa} (el que pagues encara que no consumeixis res) i el pendent $m$ és
el \emph{preu per unitat} (el que s'afegeix per cada unitat consumida).
\end{idea}

\subsection{Exemples}

\begin{exemple}[la tarifa d'un servei: quota fixa més consum]
Un servei d'atenció domiciliària cobra una \textbf{quota fixa} de $30\eur$ al mes,
més $5\eur$ per cada hora de visita. El cost mensual segons les hores $x$ és
\[
C(x)=5x+30.
\]
Aquí $m=5$ (cada hora afegeix $5\eur$) i $n=30$ (la quota fixa, el que es paga amb
$0$ hores). La gràfica és una recta que surt de $(0,30)$ i puja $5$ per cada hora:

\begin{center}
\begin{tikzpicture}
\begin{axis}[udaxis, width=9cm, height=5.4cm,
    xlabel={\small hores ($x$)}, ylabel={\small cost (\euro)},
    xmin=-0.3, xmax=9, ymin=0, ymax=75,
    xtick={0,2,4,6,8}, ytick={0,30,50,70}]
  \addplot[udblue, domain=0:8, samples=2]{5*x+30};
  \addplot[udnavy, only marks, mark=*, mark size=1.6pt] coordinates {(0,30)};
  \node[udnavy, font=\scriptsize, anchor=south east] at (axis cs:0,30) {$(0,\,30)$};
\end{axis}
\end{tikzpicture}

\figcap{La quota fixa $n=30$ és on la recta talla l'eix vertical; el pendent $m=5$ n'és la inclinació.}
\end{center}
\end{exemple}

\begin{exemple}[el signe del pendent: una prestació que baixa]
Una prestació es redueix a mesura que augmenten els ingressos: $B(x)=-20x+600$.
Ara $m=-20$ és \textbf{negatiu}, així que la recta \textbf{decreix}: per cada miler
d'euros d'ingressos més, l'ajut baixa $20\eur$. El terme $n=600$ continua sent el
punt de partida: amb $0$ ingressos, l'ajut és de $600\eur$.
\end{exemple}

\subsection{Exercicis resolts}

\begin{exresolt}[identificar $m$ i $n$ i interpretar-los]
Un casal de joves té un cost mensual $C(x)=8x+120$, on $x$ és el nombre d'inscrits.
Identifica $m$ i $n$, digues què signifiquen i calcula el cost amb $25$ inscrits.
\solucio
Comparant $C(x)=8x+120$ amb $y=mx+n$:
\[
m=8, \qquad n=120.
\]
\textbf{Interpretació:} $n=120\eur$ és el cost fix mensual (lloguer, subministraments…),
el que es paga amb $0$ inscrits. $m=8$ vol dir que cada inscrit nou afegeix $8\eur$
al cost. Per a $x=25$:
\[
C(25)=8\per 25+120=200+120=320.
\]
\resposta{$m=8$ (cost per inscrit), $n=120\eur$ (cost fix). Amb $25$ inscrits, $C=320\eur$.}
\end{exresolt}

\begin{exresolt}[escriure l'equació a partir d'una taula]
Una entitat anota els diners recaptats en una rifa solidària segons els bitllets
venuts:
\begin{center}
\begin{tabular}{c c c c}
\toprule
Bitllets $x$ & $0$ & $10$ & $20$ \\
\midrule
Recaptació $y$ ($\eur$) & $50$ & $90$ & $130$ \\
\bottomrule
\end{tabular}
\end{center}
Troba l'equació $y=mx+n$.
\solucio
El \textbf{punt de partida} (quan $x=0$) és $y=50$, per tant $n=50$. Per trobar el
\textbf{pendent}, mirem quant puja $y$ quan $x$ augmenta: de $x=0$ a $x=10$, la
recaptació passa de $50$ a $90$, és a dir, puja $40\eur$ en $10$ bitllets:
\[
m=\frac{40}{10}=4.
\]
Cada bitllet aporta $4\eur$ i hi havia una donació inicial de $50\eur$:
\[
y=4x+50.
\]
Comprovació amb $x=20$: $y=4\per 20+50=130$. Coincideix amb la taula.
\resposta{$y=4x+50$ (cada bitllet, $4\eur$; donació inicial, $50\eur$).}
\end{exresolt}

\subsection{Exercicis per fer}

\begin{exercicis}
\begin{exlist}
  \item Per a cada funció lineal, indica el pendent $m$, el terme independent $n$ i
        si creix o decreix:
    \begin{enumerate}[label=\alph*), leftmargin=1.6em, topsep=2pt]
      \item $y=3x+2$
      \item $y=-x+7$
      \item $y=0,5x$
    \end{enumerate}
  \item Una tarifa de transport adaptat costa $2\eur$ per viatge més una quota fixa
        de $15\eur$ al mes. Escriu la funció $C(x)$ del cost mensual segons els
        viatges $x$ i digues qui és $m$ i qui és $n$.
  \item Observa les dues rectes de la gràfica. Per a cadascuna, indica el terme
        independent $n$ (on talla l'eix vertical) i si el pendent és positiu o negatiu.

        \begin{center}
        \begin{tikzpicture}
        \begin{axis}[udaxis, width=9cm, height=5.4cm,
            xlabel={\small $x$}, ylabel={\small $y$},
            xmin=-0.3, xmax=9, ymin=0, ymax=9,
            xtick={0,2,4,6,8}, ytick={0,2,4,6,8}]
          \addplot[udblue, domain=0:8, samples=2]{0.7*x+1};
          \addplot[udnavy, domain=0:8, samples=2]{-0.6*x+8};
          \node[udblue, font=\scriptsize, anchor=north west] at (axis cs:6,5.2) {recta A};
          \node[udnavy, font=\scriptsize, anchor=south west] at (axis cs:5.6,4.6) {recta B};
        \end{axis}
        \end{tikzpicture}

        \figcap{Recta A i recta B: llegeix $n$ i el signe del pendent de cadascuna.}
        \end{center}

  \item Completa la taula d'una funció lineal sabent que $y=6x+10$:
        \begin{center}
        \begin{tabular}{c c c c c}
        \toprule
        $x$ & $0$ & $1$ & $2$ & $5$ \\
        \midrule
        $y$ & \rule{1.2cm}{0.4pt} & \rule{1.2cm}{0.4pt} & \rule{1.2cm}{0.4pt} & \rule{1.2cm}{0.4pt} \\
        \bottomrule
        \end{tabular}
        \end{center}
  \item \textbf{(Interpretació)} Dues tarifes de telèfon social: la tarifa P és
        $C(x)=0,10x+5$ i la tarifa Q és $C(x)=0,20x+2$, on $x$ són els minuts. Quina
        té la quota fixa més alta? Quina cobra més car cada minut? Explica què
        signifiquen $m$ i $n$ en cada cas.
  \item \textbf{(Raonament)} Una recta passa pels punts $(0,4)$ i $(2,10)$. Troba
        l'equació $y=mx+n$. (Pista: $n$ és el valor amb $x=0$; per a $m$, mira quant
        puja $y$ quan $x$ passa de $0$ a $2$.)
\end{exlist}
\end{exercicis}

% ---------------------------------------------------------------------
\fullsolucions{Sessió 4}
\begin{solucions}
\begin{sollist}
  \item (a) $m=3$, $n=2$, creix. \quad (b) $m=-1$, $n=7$, decreix. \quad
        (c) $m=0,5$, $n=0$, creix (passa per l'origen).
  \item $C(x)=2x+15$; $m=2$ (preu per viatge), $n=15\eur$ (quota fixa mensual).
  \item Recta A: $n=1$, pendent positiu (creix). \quad
        Recta B: $n=8$, pendent negatiu (decreix).
  \item $y$: \quad $10,\ 16,\ 22,\ 40$ \quad (per a $x=0,1,2,5$).
  \item La tarifa P té la quota fixa més alta ($n=5$ contra $n=2$). La tarifa Q
        cobra més car cada minut ($m=0,20$ contra $m=0,10$). $n$ és el que es paga
        sense consumir res; $m$ és el preu de cada minut.
  \item $n=4$ (valor amb $x=0$); de $x=0$ a $x=2$, $y$ puja de $4$ a $10$, és a dir
        $6$ en $2$ passos: $m=\dfrac{6}{2}=3$. Equació: $y=3x+4$.
\end{sollist}
\end{solucions}
